\chapter{Customization}\label{ch:customize}

Two ways to change things:

\begin{itemize}
  \item \textbf{Edit the repository} for changes on every machine. Configs are symlinked, so editing
        \file{~/.config/nvim/...} edits \file{~/dotfiles/...}; commit and push as usual.
  \item \textbf{Local files} for one machine only. They are not in git:
\end{itemize}

\begin{keys}{5.6cm}
\head{File}{Read by}
\file{~/.zshrc.local} & zsh, last\\
\file{~/.config/tmux/tmux.local.conf} & tmux, after the main config\\
\file{~/.config/git/config.local} & git, included at the end\\
\file{~/.config/nvim/lua/config/local.lua} & Neovim, after everything else\\
\file{~/.ssh/config.d/*} & ssh\\
\file{~/.config/apps} & \cmd{app}\\
\end{keys}

\section{Shell}

\begin{shell}
# ~/.zshrc.local
# folders searched by the sessionizer (C-a f)
export TMUX_SESSIONIZER_PATHS="$HOME/apps $HOME/web $HOME/ros2_ws"
# another tunnel config for tun
export CLOUDFLARED_CONFIG=~/.cloudflared/home.yml
# containers: don't source ROS automatically
export ROS_AUTO=0
# robots without internet: no automatic downloads in Neovim
export NVIM_OFFLINE=1
\end{shell}

\begin{itemize}
  \item Shared aliases go into \file{zsh/.config/zsh/40-aliases.zsh}.
  \item More zsh plugins: add \texttt{plug <github-user>/<repo>} to \file{20-plugins.zsh}.
  \item The prompt is \file{starship/.config/starship.toml}; add \texttt{\$hostname} to its
        \texttt{format} to show the host again.
\end{itemize}

\section{tmux}

\begin{plain}
# ~/.config/tmux/tmux.local.conf
set -g prefix C-b                 # back to the default prefix
bind C-b send-prefix
set -g status-position bottom     # status bar at the bottom
\end{plain}

\textbf{Status bar segments.} \texttt{status-right} in \file{tmux.conf} calls
\file{scripts/status.sh} once per segment: \texttt{tunnel}, \texttt{cpu}, \texttt{gpu},
\texttt{mem}, \texttt{temp}, \texttt{netspeed}, \texttt{docker}, \texttt{battery} (and \texttt{net},
the local IP, not shown by default). To add one, write a function and a line for it in the
\texttt{case} at the bottom of \file{status.sh}:

\begin{shell}
pm2_seg() {   # online PM2 processes
  command -v pm2 >/dev/null || return 0
  printf 'pm2 %s ' "$(pm2 jlist | jq '[.[] | select(.pm2_env.status == "online")] | length')"
}
case "$1" in
  pm2) pm2_seg ;;
  # ...
\end{shell}

Then add \texttt{\#(\$S/status.sh pm2)} to \texttt{status-right} and reload with \key{C-a r}.

\section{Neovim}

\textbf{Add a language} in \file{lua/config/languages.lua}: the parser in \texttt{M.treesitter},
the server in \texttt{M.servers()}, the formatter in \texttt{M.formatters\_by\_ft}. Restart Neovim;
Mason installs the server.

\textbf{Add a plugin} with a new file in \file{lua/plugins/} that returns a lazy.nvim spec;
\cmd{:Lazy} installs it and \file{lazy-lock.json} records the version.

\begin{keys}{4.4cm}
\head{What}{Where}
General keymaps & \file{lua/config/keymaps.lua}\\
Keymaps of a plugin & The \texttt{keys} of its spec in \file{lua/plugins/}\\
LSP keymaps & \file{lua/plugins/lsp.lua}\\
ROS keymaps and commands & \file{lua/util/ros.lua}\\
Options & \file{lua/config/options.lua}\\
Theme & \file{lua/plugins/colorscheme.lua} (latte, frappe, macchiato, mocha)\\
\end{keys}

\begin{plain}
-- ~/.config/nvim/lua/config/local.lua  (this machine only)
vim.g.autoformat = false
vim.opt.relativenumber = false
\end{plain}

\section{SSH hosts and git identity}

Any file in \file{~/.ssh/config.d/} is read, so hosts can be split by project; every host appears
in \cmd{s} and \key{C-a S}.

\begin{plain}
# ~/.ssh/config.d/robots
Host go2
  HostName 192.168.123.18
  User unitree

Host go2-jetson
  HostName 192.168.123.15
  User unitree
  ProxyJump go2
\end{plain}

A different git identity for work repositories:

\begin{plain}
# ~/.config/git/config.local
[includeIf "gitdir:~/ezra/"]
    path = ~/.config/git/work
\end{plain}
